% Part: sets-functions-relations
% Chapter: functions
% Section: partial-functions

\documentclass[../../../include/open-logic-section]{subfiles}

\begin{document}

\olfileid{sfr}{fun}{par}

\olsection{આંશિક વિધેયો}

\begin{explain}
ક્યારેક વિધેયની વ્યાખ્યામાં છૂટ આપવી ઉપયોગી બને છે,
જેથી દરેક સંભવિત આગત માટે વિધેયનું નિર્ગત વ્યાખ્યાયિત
હોય તે જરૂરી ન રહે. આવાં નિરૂપણોને
\emph{આંશિક વિધેયો} કહે છે.
\end{explain}

\begin{defn}
\emph{આંશિક વિધેય} $f \colon A \pto B$ એવું નિરૂપણ છે
જે $A$ના દરેક !!{element}ને $B$નો વધુમાં વધુ એક
!!{element} ફાળવે છે. જો $f$ એ $x \in A$ને $B$નો
કોઈ ઘટક ફાળવે, તો $f(x)$ \emph{વ્યાખ્યાયિત} છે
એમ કહીએ; અન્યથા તે \emph{અવ્યાખ્યાયિત} છે.
$f(x)$ વ્યાખ્યાયિત હોય તો $f(x) \fdefined$ લખીએ,
અન્યથા $f(x) \fundefined$ લખીએ. આંશિક વિધેય~$f$નો
\emph{પ્રદેશ} એ $A$નો એવો ઉપગણ છે જ્યાં તે
વ્યાખ્યાયિત હોય, એટલે કે
$\dom{f} = \Setabs{x \in A}{f(x) \fdefined}$.
\end{defn}

\begin{ex}
દરેક વિધેય $f\colon A \to B$ આંશિક વિધેય પણ છે.
$A$ પર સર્વત્ર વ્યાખ્યાયિત આંશિક વિધેયો---એટલે કે
અત્યાર સુધી આપણે જેને માત્ર વિધેય કહેતા આવ્યા છીએ તે---
\emph{સર્વત્ર વ્યાખ્યાયિત} વિધેયો પણ કહેવાય છે.
\end{ex}

\begin{ex}
$f(x) = 1/x$ દ્વારા આપેલું આંશિક વિધેય
$f \colon \Real \pto \Real$ એ $x = 0$ માટે
અવ્યાખ્યાયિત છે અને બાકીની બધી જગ્યાએ વ્યાખ્યાયિત છે.
\end{ex}

\begin{prob}
આપેલા $f\colon A \pto B$ માટે આંશિક વિધેય
$g\colon B \pto
A$ આ રીતે વ્યાખ્યાયિત કરો: કોઈ પણ $y \in B$ માટે,
જો $f(x) = y$ થાય એવો અનન્ય $x \in A$ હોય,
તો $g(y) = x$; અન્યથા $g(y) \fundefined$.
બતાવો કે જો $f$ એક-એક હોય, તો બધા $x \in \dom{f}$
માટે $g(f(x)) = x$ અને બધા $y \in \ran{f}$ માટે
$f(g(y)) = y$ થાય છે.
\end{prob}

\begin{defn}[આંશિક વિધેયનો આલેખ]
ધારો કે $f\colon A \pto B$ આંશિક વિધેય છે.
$f$નો \emph{આલેખ} એ નીચે મુજબ વ્યાખ્યાયિત સંબંધ
$R_f \subseteq A \times B$ છે:
\[
R_f = \Setabs{\tuple{x,y}}{f(x) = y}.
\]
\end{defn}

\begin{prop}
ધારો કે $R \subseteq A \times B$નો ગુણધર્મ એવો છે કે
જ્યારે પણ $Rxy$ અને $Rxy'$ હોય, ત્યારે $y = y'$ થાય છે.
તો $R$ એ આ રીતે વ્યાખ્યાયિત આંશિક વિધેય
$f\colon A \pto B$નો આલેખ છે: જો $Rxy$ થાય એવો $y$
હોય, તો $f(x) = y$; અન્યથા $f(x) \fundefined$.
જો $R$ \emph{સર્વાગત} પણ હોય, એટલે કે દરેક $x \in A$
માટે $Rxy$ થાય એવો કોઈ $y \in B$ હોય, તો $f$
સર્વત્ર વ્યાખ્યાયિત છે.
\end{prop}

\begin{proof}
ધારો કે $Rxy$ થાય એવો $y$ છે. જો $Rxy'$ થાય એવો
બીજો $y'
\neq y$ હોત, તો $R$ પરની શરતનો ભંગ થાત.
આથી, $Rxy$ થાય એવો $y$ હોય, તો તે $y$ અનન્ય છે;
તેથી $f$ સુવ્યાખ્યાયિત છે. સ્પષ્ટપણે $R_f = R$
અને જો $R$ સર્વાગત હોય, તો $f$ સર્વત્ર વ્યાખ્યાયિત છે.
\end{proof}

\end{document}
